\chapter{Συμπεράσματα \& Μελλοντική Έρευνα}
Κατα την παρούσα εργασία μελετήσαμε το πεδίο των συστημάτων εντοπισμού εσωτερικού και εξωτερικού
χώρου. Παρουσιάσαμε την προηγούμενη γνώση του τομέα αυτού στο τεχνολογικό αλλά και στο μαθηματικό
επίπεδο. Τέλος βασιζόμενοι σε αυτές τις μαθηματικές αρχές σχεδιάσαμε ένα σύστημα εντοπισμού
εσωτερικού χώρου και αξιολογήσαμε την απόδοση του.

Πιο συγκεκριμένα μιλήσαμε για τα παρακάτω θέματα
\begin{itemize}
    \item Συστήματα εντοπισμού και γιατί είναι χρήσιμα
    \item Αρχές λειτουργίας συστημάτων εντοπισμού, υπολογισμοί, θόρυβος, φιλτράρισμα και παραδείγματα
    \item Το δικό μας σύστημα και η μορφή των πειραμάτων μας
    \item Τα αποτελέσματα και η αξιολόγηση του δικού μας συστήματος
\end{itemize}

\section{Τελική Αξιολόγηση}
Πατώντας στα αποτελέσματα του προηγούμενου κεφαλαίου μπορούμε να βγάλουμε μερικά
τελικά συμπεράσματα.

\begin{itemize}
\item Απο τις δυο μεθόδους που χρησιμοποιήσαμε είναι προφανές πως η μέθοδος των
ελάχιστων τετραγώνων είναι η καλύτερη και έχει χαμηλότερη ευαισθησία στους διάφορους παράγοντες
των πειραμάτων.
\item Τα αποτελέσματα γενικά έχουν μεγάλη ευαισθησία σε περιβαλλοντικούς παράγοντες που χρειάζονται αρκετά περισσότερη μετρητική διαδικασία
για να περιγραφθούν και να ελεγχθούν
\item Σε στατικά πειράματα το σύστημα πετυχαίνει καλά αποτελέσματα όταν το \te{line of sight} είναι απρόσκοπτο (ακρίβεια περίπου μισού μέτρου)
\item Σε στατικά πειράματα το σύστημα μας δεν καταφέρνει να πετύχει ιδιαίτερα μεγάλη ακρίβεια όταν υπάρχει εμπόδιο στο
\te{line of sight} καθώς οι υπολογισμοί είναι αρκετά ευαίσθητοι στον παραπάνω θόρυβο που εισάγεται στις πρωτογενείς μεταβλητές
\item Σε πειράματα κίνησης, και στην περίπτωση με εμπόδιο και στην περίπτωση χωρίς εμπόδιο, τα αποτελέσματα δεν ειναι ιδιαίτερα ακριβή
αλλά η κίνηση του χρήστη μπορεί να διακριθεί
\end{itemize}

\section{Μελλοντική Έρευνα}
Σαν επίλογο θα προτείνουμε επιγραμματικά μερικά σημεία που θεωρούμε πως μπορούν να ερευνηθούν
στο μέλλον.
\begin{itemize}
\item Πιο ανεπτυγμένα φιλτραρίσματα, όπως το φίλτρο \te{Kalman}
\item Μια μέθοδο που αξιολογέι την ποιότητα του κάθε \te{anchor} σε πραγματικό χρόνο, ώστε το σύστημα
να χρησιμοποεί μόνο \te{anchors} που έχουν καλή ποιότητα σήματος για τον υπολογισμό της θέσης
\item Μελέτη της πειραματικής διάταξης και του υλικού ώστε τα αρχικά δεδομένα να είναι υψηλότερης
ποιότητας
\item Μελέτη μιας πειραματικής διάταξης με περισσότερα απο τέσσερα \te{anchor}
\item Συνδιασμός της προκύπτουσας πληροφορίας με εναλλακτικούς αισθητήρες, όπως επιταχυνσιόμετρα ή επιπλέον στάδια
επεξεργασίας για πρόβλεψη κίνησης
\end{itemize}
